%% kapa: exam=1 round=28 year=2017 period=1 form=A subject=civil_law subject_ko=민법 seq=32 number=32
%% source: Q-Net 감정평가사 기출문제 (cst00309) artlSeq=5199603 file=extracted/1008/5199603/2017년 제28회 감정평가사 1차시험 문제지/1교시 A형.pdf
%% conversion: kordoc 4.24.51 -> tools/gen1.py (auto); review=proofread-sonnet; flags=none
\begin{question}{32}
甲, 乙, 丙은 A토지를 각각 5분의 3, 5분의 1, 5분의 1의 지분으로 공유하고 있다. 이에 관한 설명으로 옳지 않은 것은? (다툼이 있으면 판례에 따름)
\begin{choices}
  \item 甲은 다른 공유자들과의 협의 없이 A토지의 관리방법을 정할 수 있다.
  \item 乙은 A토지에 제3자 명의로 경료된 원인무효인 근저당권설정등기의 말소를 청구할 수 있다.
  \item 등기부상의 지분과 실제의 지분이 다르고 새로운 이해관계를 가진 제3자가 없다면, 공유물분할소송에서 甲, 乙, 丙은 특별한 사정이 없는 한 실제의 지분에 따라 A토지를 분할하여야 한다.
  \item 丙의 지분 위에 근저당권이 설정된 후 A토지가 지분에 따라 분할된 때에는 특별한 합의가 없는 한 그 근저당권은 丙에게 분할된 부분에 집중된다.
  \item 甲, 乙, 丙 사이의 관리방법에 관한 약정에 따라 乙이 A토지의 특정부분만을 사용할 수 있는 경우, 특별한 사정이 없는 한 乙의 지분을 양수한 丁도 그 특정부분만을 사용할 수 있다.
\end{choices}
\end{question}
\answer{4}
